\documentclass{article}
\usepackage{graphicx} % Required for inserting images

\title{Ensayo-1 - Programación}
\author{katy Ixchely Yaxón Saloj}
\date{Febrero 2026}

\begin{document}

\maketitle
Mi área de interés para investigación en Física se centra en el estudio de los Agujeros negros y la relatividad, su estudio combina la física teórica(singularidades, horizonte de sucesos) con la observación astronómica de rayos X y ondas gravitacionales, describiendo objetos de alta masa que curvan el espacio-tiempo de manera extrema.\\


El estudio de los agujeros negros resulta especialmente relevante porque representa uno de los escenarios más extremos del universo, donde los efectos relativistas son dominantes. Fenómenos como el horizonte de eventos, la dilatación temporal gravitacional, las trayectorias de la luz (geodésicas) y la emisión de ondas gravitacionales solo pueden describirse adecuadamente mediante modelos matemáticos avanzados basados en ecuaciones diferenciales no lineales. En la mayoría de los casos, estas ecuaciones no admiten soluciones analíticas simples, lo que hace necesario el uso de métodos numéricos.\\

En este contexto, la programación se convierte en una herramienta fundamental para la investigación. Mediante el uso de lenguajes de programación y algoritmos computacionales es posible resolver numéricamente las ecuaciones de la relatividad general, simular el comportamiento de partículas y campos cerca de un agujero negro, y modelar procesos como la creación de materia o la fusión de objetos compactos. Además, la programación permite analizar datos observacionales reales, como señales de ondas gravitacionales, y compararlos con las predicciones teóricas.\\

Finalmente, la programación facilita la visualización de conceptos abstractos, como la curvatura del espacio-tiempo, lo que contribuye tanto a una mejor comprensión física como a la comunicación de resultados científicos. Por estas razones, la combinación de relatividad, agujeros negros y programación constituye un campo de investigación atractivo, actual y esencial dentro de la física moderna.\\


\end{document}
